\input amstex
\documentstyle{amsppt}
\magnification 1200
\pageno 30
\define\ogf{{G(F)}}
\define\ozf{Z_G(F)}
\define\oglf{{GL(2,E)}}
\define\oa{{A_G(F)}}
\define\oslf{{H(E)}}
\define\slf{{H(F)}}
\define\eps{\epsilon}
\define\ooo{{\Cal O}}
\define\off{{\Cal O}_F}
\define\offp{{\Cal O}_{F'}}
\define\qqq{\Bbb Q}
\define\ppp{\Bbb P}
\define\ccc{\Bbb C}
\define\rrr{\Bbb R}
\define\aaa{\Bbb A}
\define\nnn{\Bbb N}
\define\unt{\underline{T}}
\define\unx{\underline{x}}
\define\oga{{G(\aaa )}}
\define\oba{{B_G(\aaa )}}
\define\oza{{Z_G(\aaa )}}
\define\ogla{{GL(2,\aaa_L )}}
\define\osla{{H(\aaa_L )}}
\define\sla{H(\aaa )}
\define\oee{{\Cal O}_{E}}
\define\oeep{{\Cal O}_{E'}}
\define\VV{\Cal V}
\define\HH{\Cal H}
\hsize 6.5 true in
\vsize 8.7 true in
\vskip .2in
\subheading{\S 4 Regularity of patches for $G=G_2$}
\vskip .1in
It is not difficult to write all these patches of \S 3
down and show that they are
regular for $G_2$. Recall $\Sigma_G^+=\{\alpha,\beta,
\gamma:=\alpha+\beta,\delta:=2\alpha +\beta,
\epsilon:=3\alpha +\beta,\zeta:=3\alpha +2\beta\}$, with
$\alpha$ short and $\beta$ long, so $p=6$ and dim $Z=7$. We have variables
$$
\split
&x(\alpha),x(\beta),
x(\gamma),x(\delta),x(\epsilon),x(\zeta)\\
&z(\alpha),z(\beta),
w(\gamma),w(\delta),w(\epsilon),w(\zeta)\\
&\lambda,
\endsplit
\tag 4.1
$$
subject to the usual relations (3.1).
\par
With $B_+$ choosen as in \S 3, we define $\omega_{Y''(B_+)}$ to be the
$G$-invariant form
$$
\omega_{Y''(B_+)}=d\lambda \wedge
\wedge
dx(\alpha)\wedge dx(\beta)
\wedge dx(\gamma)\wedge dx(\delta)\wedge dx(\epsilon)\wedge dx(\zeta)
\wedge dv,
$$
where $dv$ is a form on $N_+$. We use this to compute the Igusa constants
$\beta$ in this section as in \cite{H2}.
\par
As we know (from lemma 3.5), the Richardson patches $Z_{S_1}$ are regular.
Now let us consider the remaining patches. Since $S_1\cup S_2\not= \Delta$
if $S_2\not= \emptyset$, it is easy to list
them all: $S_1=\emptyset$, $S_2=\{\alpha\}$, where
$\sigma(\alpha)=\alpha,\gamma,\delta,\epsilon,\zeta$, or
$S_1=\emptyset$, $S_2=\{\beta\}$, where
$\sigma(\beta)=\beta,\gamma,\delta,\epsilon,\zeta$.
This may be deduced from the discussion above. We next list the generators of
the local coordinate ring in each case:
\newline
$S_2=\{\alpha\}, \sigma(\alpha)=\alpha$: Same as $Z_{\{\alpha\}}$. The divisor
$\lambda =0$ occurs with multiplicity one.
\newline
$S_2=\{\alpha\}, \sigma(\alpha)=\gamma$: Generators: $z(\alpha), z(\beta),
x(\gamma), w(\gamma), x(\delta), x(\epsilon),
x(\zeta)$. Regular at $\lambda =0$. The divisor $\lambda =0$
occurs with multiplicity one (represented by $z(\alpha)=0$).
\newline
$S_2=\{\alpha\}, \sigma(\alpha)=\delta$: Generators: $z(\alpha), z(\beta),
w(\gamma), x(\delta),w(\delta), x(\epsilon),
x(\zeta)$. Regular at $\lambda =0$. The divisor $\lambda =0$
occurs (as $z(\alpha)=0$) with multiplicity two.
\newline
$S_2=\{\alpha\}, \sigma(\alpha)=\epsilon$: Generators: $z(\alpha), z(\beta),
w(\gamma), w(\delta),w(\epsilon),x(\epsilon),x(\zeta)$.
Regular at $\lambda =0$. The divisor $\lambda =0$
occurs (as $z(\alpha)=0$) with multiplicity three.
\newline
$S_2=\{\alpha\}, \sigma(\alpha)=\zeta$: Generators: $z(\alpha), z(\beta),
w(\gamma), w(\delta),w(\epsilon), w(\zeta),
x(\zeta)$. Regular at $\lambda =0$. The divisor $\lambda =0$
occurs (as $z(\alpha)=0$) with multiplicity three.
\newline
$S_2=\{\beta\}, \sigma(\beta)=\beta$: Same as $Z_{\{\beta\}}$.
The divisor $\lambda =0$
occurs with multiplicity one.
\newline
$S_2=\{\beta\}, \sigma(\beta)=\gamma$: Generators: $z(\alpha), z(\beta),
x(\gamma), w(\gamma), x(\delta), x(\epsilon),
x(\zeta)$. Regular at $\lambda =0$. The divisor $\lambda =0$
occurs with multiplicity one (represented by $z(\beta)=0$).
\newline
$S_2=\{\beta\}, \sigma(\beta)=\delta$: Generators: $z(\alpha), z(\beta),
w(\gamma), x(\delta),w(\delta), x(\epsilon),
x(\zeta)$. Regular at $\lambda =0$. The divisor $\lambda =0$
occurs (as $z(\beta)=0$) with multiplicity one.
\newline
$S_2=\{\beta\}, \sigma(\beta)=\epsilon$: Generators: $z(\alpha), z(\beta),
w(\gamma), w(\delta),w(\epsilon),x(\epsilon),x(\zeta)$.
Regular at $\lambda =0$. The divisor $\lambda =0$
occurs (as $z(\beta)=0$) with multiplicity one.
\newline
$S_2=\{\beta\}, \sigma(\beta)=\zeta$: Generators: $z(\alpha), z(\beta),
w(\gamma), w(\delta),w(\epsilon), w(\zeta),
x(\zeta)$. Regular at $\lambda =0$. The divisor $\lambda =0$
occurs (as $z(\beta)=0$) with multiplicity two.
\par
There is some redundancy here since $E_{\emptyset,\{\beta\},\sigma}=
E_{\{\beta\}}$ if $\sigma(\beta)=\beta,\gamma,\delta,\epsilon$.
\par
We denote the (irreducible) branch of $\{\lambda =0\}\subset Z$
occuring in $Z_{S_1,S_2,\sigma}$ by $E_{S_1,S_2,\sigma}$. The same notation
will be used for the corresponding branch in $Y_U$ and for its Zariski closure
in $Y_1$. To determine which of the $\kappa$-stable germs, if any,
the divisor $E_{S_1,S_2,\sigma}$
contributes to by formula 0.6,
we must compute the Igusa constant
$$
\beta(\tilde{E}_{S_1,S_2,\sigma})=
{b(\tilde{E}_{S_1,S_2,\sigma})\over a(\tilde{E}_{S_1,S_2,\sigma})}
\tag 4.2
$$
where $\tilde{E}$ denotes the preimage of $E$ in the
desingularized variety $Y_C$. 

We record two observations regarding the Igusa constants $\beta(E)$.
\newline
(I)
First, if $\phi:V\rightarrow W$ is an isomorphism of smooth varieties
which sends a divisor $E\subset V$ to a divisor $F=\phi(E)\subset W$
then $\beta(E)=\beta(F)$.
\newline
(II)
Secondly, let $V'$ denote a blow up of a variety
$V$ with exceptional divisor $E_{V'}$.
Recall that for each irreducible branch $E$ on $V$,
the pull-back ${\xi'}^*$ of the blow down map $\xi':V'\rightarrow V$
satisfies ${\xi'}^*(E)=\tilde{E}+\delta E_{V'}$, where $\tilde{E}$
is the strict transform
of $E$, and
$$
\delta:=
\cases 1, &\text{ if $E$ contains the center of the blow up},\\
0, &\text{ otherwise}.
\endcases
$$
A repeated use of this global identity can be used to obtain
local information about the
Igusa constants on $Z$.

For example, if $E$ denotes a branch of $\{\lambda =0\}\subset Z$ and
$\tilde{E}$ its strict transform in a desingularization $\tilde{Z}$,
we have
$$
\beta(\tilde{E})=\beta(E),
$$
provided $E$ does not contain any of the images of the centers of
the blow ups used to create $\tilde{Z}$. The only other divisors
in $\{\lambda =0\}\subset \tilde{Z}$ other than the strict
transforms $\tilde{E}$ above are the exceptional divisors.
To compute all the $\beta$-values of the
branches of $\phi^{-1}(1)$ in $Y_C$ we need only compute the $\beta$-values
of the divisors $E_{S_1,S_2,\sigma}$ on $Z$ and the branches of
the exceptional divisors of the blow-up.

We digress to
compute $\beta({E}_{S_1,S_2,\sigma})$ on $Z_{S_1,S_2,\sigma}$. We
need not compute the $\beta$-values for the exceptional divisors
because of lemma 0.8 and the fact that all the needed $F$-rational
${\Cal O}$-divisors can be accounted for among the $E_{S_1,S_2,\sigma}$.

If $S_2=\emptyset$ then $\beta({E}_{S_1,S_2,\sigma})$
has been computed in \cite{H1, ch II, Proposition 9.1}
so some of our calculations are covered by Hales' general
calculations. For completeness, we list everything anyway.
\vskip .1in
{\bf Igusa forms on the fundamental patches of $Y''$}
\vskip .2in
We pull-back the patch $Z_{S_1,S_2,\sigma}$ 
to $Y''(B_+,B_-)$ via the open immersion
$Y''(B_+,B_-)\rightarrow Z$. Call this inverse image $Y''_{S_1,S_2,\sigma}$.
The Igusa form $\omega_U$ on $U=Y''_{S_1,S_2,\sigma}$ is related to the
Igusa form $\omega_{Z,U}$ on $Z_{S_1,S_2,\sigma}$ by 
$\omega_U=\omega_{Z,U}\wedge dv$,
where $dv$ is regular. We compute $\omega_{Z,U}$below.

$Z_{\{\alpha,\beta\}}$: Here the multiplicity $a(E_{\{\alpha,\beta\}})=1$
and $b(E_{\{\alpha,\beta\}})=1$ since
$$
\omega_{Z,U}
=d\lambda\wedge dx(\alpha)\wedge dx(\beta)\wedge dx(\gamma)
\wedge dx(\delta)\wedge dx(\epsilon)\wedge dx(\zeta) .
$$
Therefore, $\beta(E_{\{\alpha,\beta\}})=1$.
\vskip .1in
$Z_{\alpha}$: Here the multiplicity $a(E_{\alpha})=1$
and $b(E_{\alpha})=2$ since
$$
\split
\omega_{Z,U}
&=d\lambda\wedge dx(\alpha)\wedge dx(\beta)\wedge dx(\gamma)
\wedge dx(\delta)\wedge dx(\epsilon)\wedge dx(\zeta)\\
&=-{\lambda\over z(\beta)^2}
d\lambda\wedge dx(\alpha)\wedge dz(\beta)\wedge dx(\gamma)
\wedge dx(\delta)\wedge dx(\epsilon)\wedge dx(\zeta).
\endsplit
$$
Therefore, $\beta(E_{\alpha})=2$.
\vskip .1in
$Z_{\beta}$: Here the multiplicity $a(E_{\beta})=1$
and $b(E_{\beta})=2$ since
$$
\split
\omega_{Z,U}
&=d\lambda\wedge dx(\alpha)\wedge dx(\beta)\wedge dx(\gamma)
\wedge dx(\delta)\wedge dx(\epsilon)\wedge dx(\zeta)\\
&=-{\lambda\over z(\alpha)^2}
d\lambda\wedge dz(\alpha)\wedge dx(\beta)\wedge dx(\gamma)
\wedge dx(\delta)\wedge dx(\epsilon)\wedge dx(\zeta).
\endsplit
$$
Therefore, $\beta(E_{\beta})=2$.
\vskip .1in
$Z_{\emptyset,\alpha,\sigma}, \ \sigma(\alpha)=\gamma$:
Here the multiplicity $a(E_{\emptyset,\alpha,\sigma})=1$, by the
above list. Also, $b(E_{\emptyset,\alpha,\sigma})=2$ since
$$
\split
\omega_{Z,U}
&=d\lambda\wedge dx(\alpha)\wedge dx(\beta)\wedge dx(\gamma)
\wedge dx(\delta)\wedge dx(\epsilon)\wedge dx(\zeta)\\
&=c{x(\gamma)^3z(\beta)z(\alpha)\over w(\gamma)^4}
dz(\alpha)\wedge dz(\beta)\wedge dw(\gamma)\wedge dx(\gamma)
\wedge dx(\delta)\wedge dx(\epsilon)\wedge dx(\zeta).
\endsplit
$$
Therefore, $\beta(E_{\emptyset,\alpha,\sigma})=2$.
\vskip .1in
$Z_{\emptyset,\alpha,\sigma}, \ \sigma(\alpha)=\delta$:
Here the multiplicity $a(E_{\emptyset,\alpha,\sigma})=2$, by the
above list. Also, $b(E_{\emptyset,\alpha,\sigma})=6$ since
$$
\split
\omega_{Z,U}
&=d\lambda\wedge dx(\alpha)\wedge dx(\beta)\wedge dx(\gamma)
\wedge dx(\delta)\wedge dx(\epsilon)\wedge dx(\zeta)\\
&=c{x(\delta)^4z(\beta)z(\alpha)^5\over w(\delta)^5}
dz(\alpha)\wedge dz(\beta)\wedge dw(\gamma)\wedge dw(\delta)
\wedge dx(\delta)\wedge dx(\epsilon)\wedge dx(\zeta).
\endsplit
$$
Therefore, $\beta(E_{\emptyset,\alpha,\sigma})=3$.
\vskip .1in
$Z_{\emptyset,\alpha,\sigma}, \ \sigma(\alpha)=\epsilon$:
Here the multiplicity $a(E_{\emptyset,\alpha,\sigma})=3$, by the
above list. Also, $b(E_{\emptyset,\alpha,\sigma})=10$ since
$$
\split
\omega_{Z,U}
&=d\lambda\wedge dx(\alpha)\wedge dx(\beta)\wedge dx(\gamma)
\wedge dx(\delta)\wedge dx(\epsilon)\wedge dx(\zeta)\\
&=c{x(\epsilon)^5z(\beta)z(\alpha)^9\over w(\delta)^6}
dz(\alpha)\wedge dz(\beta)\wedge dw(\gamma)\wedge dw(\delta)
\wedge dw(\epsilon)\wedge dx(\epsilon)\wedge dx(\zeta).
\endsplit
$$
Therefore, $\beta(E_{\emptyset,\alpha,\sigma})=10/3$.
\vskip .1in
$Z_{\emptyset,\alpha,\sigma}, \ \sigma(\alpha)=\zeta$:
Here the multiplicity $a(E_{\emptyset,\alpha,\sigma})=3$, by the
above list. Also, $b(E_{\emptyset,\alpha,\sigma})=11$ since
$$
\split
\omega_{Z,U}
&=d\lambda\wedge dx(\alpha)\wedge dx(\beta)\wedge dx(\gamma)
\wedge dx(\delta)\wedge dx(\epsilon)\wedge dx(\zeta)\\
&=c{x(\zeta)^6z(\beta)^7z(\alpha)^{10}\over w(\delta)^7}
dz(\alpha)\wedge dz(\beta)\wedge dw(\gamma)\wedge dw(\delta)
\wedge dw(\epsilon)\wedge dw(\zeta)\wedge dx(\zeta).
\endsplit
$$
Therefore, $\beta(E_{\emptyset,\alpha,\sigma})=11/3$.
\vskip .1in
$Z_{\emptyset,\beta,\sigma}, \ \sigma(\beta)=\zeta$:
Here the multiplicity $a(E_{\emptyset,\beta,\sigma})=2$, by the
above list. Also, $b(E_{\emptyset,\beta,\sigma})=8$ since
$$
\split
\omega_{Z,U}
&=d\lambda\wedge dx(\alpha)\wedge dx(\beta)\wedge dx(\gamma)
\wedge dx(\delta)\wedge dx(\epsilon)\wedge dx(\zeta)\\
&=c{x(\zeta)^6z(\beta)^7z(\alpha)^{10}\over w(\delta)^7}
dz(\alpha)\wedge dz(\beta)\wedge dw(\gamma)\wedge dw(\delta)
\wedge dw(\epsilon)\wedge dw(\zeta)\wedge dx(\zeta).
\endsplit
$$
Therefore, $\beta(E_{\emptyset,\beta,\sigma})=4$.
\vskip .1in
We summarize this in the following table:
$$
\matrix
S_1&S_2&\sigma&\beta(E_{S_1,S_2,\sigma})-1\\
 & & &\\
\alpha,\beta&\emptyset&------&0\\
\alpha&\emptyset&------&1\\
\beta&\emptyset&------&1\\
\emptyset&\emptyset&------&\text{trivial}\\
\emptyset&\alpha&\alpha&1\\
\emptyset&\alpha&\gamma&1\\
\emptyset&\alpha&\delta&2\\
\emptyset&\alpha&\epsilon&7/3\\
\emptyset&\alpha&\zeta&8/3\\
\emptyset&\beta&\beta&1\\
\emptyset&\beta&\gamma&1\\
\emptyset&\beta&\delta&1\\
\emptyset&\beta&\epsilon&1\\
\emptyset&\beta&\zeta&3.
\endmatrix
\tag 4.3
$$
As was remarked above, there is some redundancy here... . 

Let $H_1=PGL(3)$
and $H_2=(SL(2)\times SL(2))/(\pm 1)$.
\proclaim{Lemma 4.4} In the notation of Lemmas 3.5 and 3.6, we have
$$
\beta(E_S)={b(E_S)\over a(E_S)}=1+|\Sigma_S^+|=1+{ dim(Cent_G(u))-2\over 2},
$$
where $u$ is an element of the Richardson orbit $\ooo_S$ associated to the 
parabolic $P_S$ and the divisor $E_S$ is the
Richardson divisor contained in the patch $Z_S$.
The Igusa constants for 
$$
E_{\ooo}=
\cases
E_{\emptyset,\alpha,\delta}, & \ooo \text{ is $2$-regular},\\
E_{\emptyset,\beta,\zeta}, &  \ooo \text{ is $3$-regular},
\endcases
$$
are given by
$$
\beta(E_{\ooo})=
{b(E_{\ooo})\over a(E_{\ooo})}
=1+|\Sigma_{H_i}^+|=1+{ dim(Cent_G(u))-2\over 2},
$$
where $i=1,2$ is determined by the orbit as indicated below and
$u$ is an element of the induced orbit 
$$
\ooo=
\cases
Ind_{H_2}^G(1), &\ooo \text{ is $2$-regular},\\
Ind_{H_1}^G(1), &\ooo \text{ is $2$-regular},
\endcases
$$
associated to the endoscopic $H_i$ and the divisor $E_{\ooo}$.
\endproclaim
\demo{proof} The first part of the lemma is a special case of 
\cite{H1,ch II, Proposition 9.1} (it also follows
from our case-by-case calculations).
The second part is verified 
by the calculations above. (For various 
facts about unipotent orbits, including the definition of
induction used here, see \cite{Spal}.) \qed
\enddemo
\enddocument

\vskip .1in
{\bf Igusa forms on $Y_1$}
\vskip .2in
Each of the Igusa forms $\omega_U$ on $Y''(B_+,B_-)$ can be pulled-back to 
$Y_1(B_+,B_-)$ via the coordinate relations (1.22). 
This allows us to compute the dependence of $\omega_U$ on $X$. 
The result of this computation is essentially
independent of the patch as we will see.

On the preimage $U=Y_{1,S_1,S_2,\sigma}$ of $Z_{S_1,S_2,\sigma}$ under 
the open immersion $Y_{1,S_1,S_2,\sigma}\rightarrow Z$, we have
$\omega_{U}=\omega_{Z,U}\wedge dv$, where $dv$ is regular and
$$
\omega_{Z,U}=h(t_0,\lambda)d\lambda\wedge dn_\alpha\wedge dn_\beta\wedge
...\wedge dn_\zeta ,
$$
where $t_0$ is as in (1.17), $t_0=exp(\lambda X)$, and
$$
\split
h(t_0,\lambda)
&=(1-\alpha(t_0)^{-1})...(1-\zeta(t_0)^{-1})+
\text{ higher order terms in }\lambda \\
&=\lambda^6 \alpha(X)...\zeta(X)+ \text{ higher order terms in }\lambda.
\endsplit
$$
\enddocument
